\documentclass[10pt,a4paper]{amsart}
\usepackage[margin=19mm]{geometry}
\usepackage{amsmath,amssymb,mathtools}
\usepackage[scr=rsfs,cal=euler]{mathalfa}
\usepackage{xfrac}
\usepackage[unicode,hidelinks,bookmarks=false]{hyperref}
\newcommand{\ie}{i.e.\ }
\newcommand{\cf}{cf.\ }
\newcommand{\resp}{respectively\ }
\newcommand{\Subsection}{\subsection}
\providecommand{\nobreakdash}{\nobreak\hbox{-}}
\setlength{\emergencystretch}{2em}
\multlinegap0pt
\usepackage{bm}
\usepackage{bbm}
\usepackage{dsfont}
\usepackage{xspace}
\usepackage{cancel}
\usepackage{mathtools}
\usepackage{amsthm}
\makeatletter
\@ifundefined{theorem}{\newtheorem{theorem}{Theorem}}{}
\makeatother
\newcommand{\range}{\mathrm{range}}
\newcommand{\sA}{\mathcal{A}}
\title[Classification of Double Saddle-Point Systems]{Classification of Double Saddle-Point Systems}
\author{Susanne Bradley, Chen Greif}
\date{Source: arXiv:2605.14157v1 (CC BY 4.0)}
\begin{document}
\maketitle

\noindent\emph{Source and licence.} Classification of Double Saddle-Point Systems. Version arXiv:2605.14157v1 (CC BY 4.0), distributed under the Creative Commons Attribution 4.0 International licence, as recorded in the arXiv licence metadata for this exact version and shown on the abstract page https://arxiv.org/abs/2605.14157v1. Full rights notice: licenses/arxiv-2605.14157-CC-BY-4.0.txt.\par\smallskip

\noindent\textit{Editorial scope.} Counted unit: the Theorem whose source label is \texttt{None} in the article (source0.tex lines 1451--1453, source label \texttt{None}), delivered together with the article's own complete original proof of it (source0.tex lines 1455--1482, one \texttt{proof} environment of 28 lines). No mathematics is written, repaired or re-derived: statement and proof are the article's own text line for line. Required context retained uncounted: none. Every definition, display and label of those lines is retained unchanged. Omitted from this package and not redistributed: 1--1450, 1454--1454, 1483--1670. Nothing inside a retained excerpt is omitted or altered: every retained body is delivered exactly as the article's lines are written, so no omission receipt of any kind accompanies this package. Citations: the retained excerpts carry no citation command at all, so no bibliography entry of the article is transcribed and no reference is redistributed. Reference adaptation: none was needed and none was made; every reference inside the delivered unit resolves inside it, so no unresolved reference or citation is produced. Printed slips kept literal: no printed word, symbol, hypothesis or number of the counted statement or of its proof was altered. Layout and class: the article's own class is \texttt{article} with packages including lipsum, amsfonts, amsmath, amsthm, graphicx, epstopdf, algorithmic, todonotes, enumitem, comment, fancybox; the wrapper is the lane's pinned \texttt{amsart} editorial wrapper at 10pt on A4, which restates the theorem-like environments the retained text needs and adds the article's own macro definitions that the retained text needs (\texttt{\textbackslash range}, \texttt{\textbackslash sA}), with extra wrapper packages (bm, bbm, dsfont, xspace, cancel, mathtools). The delivered numbering is the unit's own. Assets: no retained span contains a figure environment, a graphics-inclusion command, a PDF-page inclusion, a table or a TikZ picture; no image file is copied into this package, the package declares no asset dependency and no image file is redistributed with it. Licence: version arXiv:2605.14157v1 of this article is distributed under the Creative Commons Attribution 4.0 International licence, as recorded in the arXiv licence metadata for this exact version and shown on the abstract page https://arxiv.org/abs/2605.14157v1. A full rights notice is delivered at licenses/arxiv-2605.14157-CC-BY-4.0.txt. Characters: every retained excerpt is pure ASCII, so the delivered engine, XeLaTeX with its default Latin Modern text font, reports no missing character.
\begin{theorem}[Eigenvalues of $(\mathcal{P}_{D,2}^a)^{-1} \sA_0$]
Assume without loss of generality that $n_2 \ge n_3$. The preconditioned matrix $(\mathcal{P}_{D,2}^a)^{-1} \sA_0$ has eigenvalue $\lambda = 1$ with multiplicity $n_1-(n_2+n_3)$, and eigenvalues $\lambda = \frac{1 \pm \sqrt{5}}{2}$, each with multiplicity of at least $n_2-n_3$.
\end{theorem}

\begin{proof}
    We begin by writing the eigenvalue equations:
    \begin{subequations}\label{eq:eig_blocktri_noreg_prec}
        \begin{align}
            A_1 x + B_1^T y + B_2^T z &= \lambda A_1 x; \label{eq:eig_a1x}\\ 
            B_1 x &= \lambda S_{a,1} y; \label{eq:eig_sa1y}\\ 
            B_2 x &= \lambda S_{a,2} z. \label{eq:eig_sa2z}
        \end{align}
    \end{subequations}
    The stated result for the eigenvalue $\lambda =1$ follows from setting $x \in \ker\left(\begin{bmatrix} B_1 \\ B_2 \end{bmatrix} \right)$, $y=0$, $z=0$.

    Next, consider $x \in \ker(B_2)$ but $x \notin \ker(B_1)$. Equation \eqref{eq:eig_sa2z} gives us $z=0$. Using \eqref{eq:eig_sa1y} to obtain $y = \frac{1}{\lambda}S_{a,1}^{-1}B_1x$ and substituting this into \eqref{eq:eig_a1x} gives us
    \[
    x + \frac{1}{\lambda} A_1^{-1}B_1^T S_{a,1}^{-1}B_1 x - \lambda x = 0.
    \]
    Observe that, when $S_{a,1} = B_1A_1^{-1}B_1^T$, the matrix $A_1^{-1}B_1^T S_{a,1}^{-1}B_1$ is a projector onto $\range(A_1^{-1}B_1^T)$.
    For $x$ in this range, the equation for $\lambda$ simplifies to
    \[
    \lambda - 1 - \frac{1}{\lambda} = 0,
    \]
    yielding the values $\lambda = \frac{1\pm \sqrt{5}}{2}$.

    The multiplicity of these eigenvalues is given by the dimension of the subspace $\range(A_1^{-1} B_1^T) \cap \ker(B_2)$ (because we required $x \in \ker(B_2)$ in this case). We observe that a vector $x$ in $\range(A_1^{-1} B_1^T)$ can be written as $x = A_1^{-1} B_1^Ty$, for some $y \in \mathbb{R}^{n_2}$. Such a vector lies in $\ker(B_2)$ if and only if
    \[
    B_2 A_1^{-1}B_1^T y = 0.
    \]
    Thus, the multiplicity of the eigenvalues is given by the nullity of the matrix \\ $B_2 A_1^{-1}B_1^T \in \mathbb{R}^{n_3 \times n_2}$, which is at least $n_2 - n_3$.
\end{proof}

\end{document}
